\section{Pole observable and electroweak scheme separation}
\label{sec:pole}
% Status: cited/open. Observable and placement convention are defined; numerical inputs require source audit.

The observable used for comparison is the pole-mass ratio
\begin{equation}
  \sPole = 1-\frac{M_{W,\rm pole}^2}{M_{Z,\rm pole}^2}.
  \label{eq:spole}
\end{equation}
This ratio is distinct from a running \(\MSbar\) weak angle and from effective weak mixing angles extracted from neutral-current asymmetries.  The manuscript therefore treats Eq.~\eqref{eq:spole} as the primary target and later relates other schemes to it.

The following placement convention will be used.  The Rivero--de Vries radical is read as a spectral object attached to the W/Z pole data.  Other readings could attach the same radical to a high-scale boundary value, to a running \(\MSbar\) coupling, or to an internal compactification eigenvalue.  Each reading changes the physical content of the proposal.  The pole reading gives the word ``spectrum'' its literal electroweak meaning.

A vector-boson pole is defined by the complex pole of the propagator,
\begin{equation}
  s_V=M_{V,\rm pole}^2-iM_{V,\rm pole}\Gamma_{V,\rm pole}.
  \label{eq:complex-pole}
\end{equation}
Throughout this manuscript, \(M_{V,\rm pole}\) and
\(\Gamma_{V,\rm pole}\) refer to the parameterization in
Eq.~\eqref{eq:complex-pole}.  Martin and Robertson also record the
alternative reporting convention
\begin{equation}
  s_{\rm pole}=\left(M-i\Gamma/2\right)^2 .
  \label{eq:alternate-pole-parameterization}
\end{equation}
Any input reported in that convention must be translated to
Eq.~\eqref{eq:complex-pole} before Eq.~\eqref{eq:spole} is formed
\cite{MartinRobertson2025}.
For an unmixed charged vector, \(s_V\) is the zero of the inverse transverse
propagator after renormalization,
\begin{equation}
  \Delta^{-1}_{T,V}(s_V)=0.
  \label{eq:pole-condition}
\end{equation}
In the neutral sector, \(s_Z\) is the massive eigenvalue pole of the transverse
\((\gamma,Z)\) inverse-propagator matrix.  Equivalently, after the photon null
direction is separated,
\begin{equation}
  \det \Delta^{-1}_{T,\gamma Z}(s_Z)=0
  \quad\hbox{on the massive neutral subspace}.
  \label{eq:neutral-pole-condition}
\end{equation}
When experimental inputs are quoted in a variable-width Breit--Wigner
line-shape convention, the conversion to pole parameters must be performed
before using Eq.~\eqref{eq:spole}.  With a common W/Z variable-width
convention, the following equations define a conditional convention map with
\(\gamma_V\) formed from the quoted variable-width parameters:
\begin{equation}
  \gamma_V=\frac{\Gamma_{V,\rm BW}}{M_{V,\rm BW}} .
  \label{eq:bw-width-ratio}
\end{equation}
The pole parameters in Eq.~\eqref{eq:complex-pole} are then obtained as
\begin{equation}
  M_{V,\rm pole}
  =
  \frac{M_{V,\rm BW}}{\sqrt{1+\gamma_V^2}},
  \qquad
  \Gamma_{V,\rm pole}
  =
  \frac{\Gamma_{V,\rm BW}}{\sqrt{1+\gamma_V^2}} .
  \label{eq:bw-pole-mass-width}
\end{equation}
Martin and Robertson summarize the corresponding vector-boson mass relation as
\begin{equation}
  M_{V,\rm BW}^2=M_{V,\rm pole}^2+\Gamma_{V,\rm pole}^2,
  \qquad V=W,Z .
  \label{eq:bw-pole-relation}
\end{equation}
Equations~\eqref{eq:bw-pole-mass-width}--\eqref{eq:bw-pole-relation}
apply after the quoted W and Z inputs have been verified to share the same
variable-width Breit--Wigner convention.  Fixed-width line-shape inputs,
already-pole inputs, and inputs using Eq.~\eqref{eq:alternate-pole-parameterization}
require their stated convention map.  This conversion is part of the source
audit for the pole comparison \cite{MartinRobertson2025}.

The resulting convention chain is
\begin{equation}
  (M_{V,\rm BW},\Gamma_{V,\rm BW})
  \longrightarrow
  (M_{V,\rm pole},\Gamma_{V,\rm pole})
  \longrightarrow
  s_V
  \longrightarrow
  \sPole .
  \label{eq:convention-chain}
\end{equation}

\subsection{Scheme map}

The comparison uses the following scheme map before any weak-angle quantity is interpreted:
\begin{center}
\small
\begin{tabular}{p{0.23\linewidth}p{0.31\linewidth}p{0.34\linewidth}}
\toprule
Object & Definition level & Role in this manuscript\\
\midrule
\(\sPole\) & complex-pole mass ratio & primary low-energy spectral target\\
Breit--Wigner masses & experimental line-shape convention & input convention requiring Eq.~\eqref{eq:bw-pole-relation}\\
\(\hat s^2(Q)_{\MSbar}\) & running coupling ratio in a renormalization scheme & scheme translation and loop-scale comparison\\
\(\sin^2\theta_{\rm eff}^{\ell}\) & neutral-current effective angle & phenomenological comparison after electroweak corrections\\
GUT boundary value & high-scale representation datum & comparison class for structural weak-angle values\\
compactification eigenvalue & internal spectral datum & possible origin for \(\mathcal O_J\) after a dynamical derivation\\
\bottomrule
\end{tabular}
\end{center}

The pole reading selects the first row.  The remaining rows are translation targets.  A future source-audited electroweak comparison should begin from the quoted experimental convention, map to complex pole quantities, form Eq.~\eqref{eq:spole}, and then compare with Eq.~\eqref{eq:sdvdef}.

The inherited numerical illustration uses
\begin{equation}
  M_W=M_Z\sqrt{1-\sdV}.
  \label{eq:MWfromMZ}
\end{equation}
The body of the paper defers a current numerical comparison until the quoted
mass convention, complex-pole conversion, and uncertainty propagation are
source-audited together.  Appendix~\ref{app:numerics} records
the inherited arithmetic for provenance.  The CDF-II measurement is relevant
only as an example showing that different line-shape inputs can move the
descriptive pole-ratio comparison \cite{CDFII2022}.

The Standard Model parameter literature distinguishes several vacuum and renormalization conventions.  In particular, pure \(\MSbar\) treatments may define the Higgs vacuum expectation value as the minimum of the loop-corrected effective potential in Landau gauge, eliminating tadpole diagrams at the cost of gauge-choice dependence in intermediate parameters \cite{MartinRobertson2025}.  Such scheme details are relevant for comparing a mass-ratio statement with Lagrangian parameters, while Eq.~\eqref{eq:spole} remains the pole-spectrum target.

Loop effects are therefore a conceptual discriminator.  Pole, effective, and running definitions are related by scheme-dependent self-energies and threshold terms.  They ask what object the radical belongs to and at what scale it is to be read.  A high-scale reading would make the DeVries value analogous to a GUT boundary datum.  The pole-data reading makes it a low-energy spectral datum and transfers the burden to string/Kaluza--Klein dynamics.

\subsection{Limited claim}

The formal claim of this section is limited:
\begin{equation}
  \sdV
  \quad\text{is compared first with}\quad
  1-\frac{M_{W,\rm pole}^2}{M_{Z,\rm pole}^2}.
  \label{eq:pole-claim-limited}
\end{equation}
Every comparison to \(\MSbar\) inputs, effective angles, or line-shape masses requires an explicit map.  Thus a disagreement among schemes diagnoses the placement of \(\sqrt{J^2+4J}\) and the scale at which the radical is read.
